% !TEX root = ../../../tree_notes.tex
% =====================================================================
%  12.1  Depth, Height and Subtree Size in One DFS  sec:c09-measures
%  Exports: def:c09-measures  thm:c09-onepass
%           lst:c09-dfs3  lst:c09-iter  fig:c09-triples
%           ex:c09-measure-run
%           prob:c09-deep  lst:c09-742-soln   (practices: LC 2331, LC 1609)
%  lst:c09-iter is the listing Chapter 16's kit cites for 5*10^5 chains
%  (forward contract discharged here). prob:c09-deep replaces the stale
%  LC 865 pointer (URL owned by Ch.2); ch11 prose was synced in the
%  same pass. Binary-tree recursions stay Ch.1's -- this section is the
%  GENERAL-tree re-framing: adjacency list + parent guard + one DFS.
% =====================================================================
\section{Depth, Height and Subtree Size --- One DFS for All Three}\label{sec:c09-measures}

\begin{previouslybox}
\begin{enumerate}[nosep, leftmargin=1.6em]
  \item Chapter~1 defined $\depth(v)$, $\height(v)$ and $\size(v)$ on rooted trees and
        proved the recurrences $\height(v)=1+\max_c\height(c)$, $\size(v)=1+\sum_c\size(c)$
        (\thmref{c01-height-recursion}). Which two conventions made a leaf need no special
        case? ($\height(\nil)=-1$, $\size(\nil)=0$)
  \item Chapter~2's bottom-up pattern lets a node read answers its children computed.
        Which of our three measures travels \emph{down} the recursion and which travel
        \emph{up}? (depth down as a parameter; height and size up as return values)
  \item A general tree arrives as an undirected adjacency list (\chref{c04}). What single
        guard turns ``walk the graph'' back into ``walk the tree''? (skip the parent)
  \item Run recursive DFS on a chain of $10^{6}$ nodes. What dies first -- the algorithm
        or the machine? (the call stack; the rewrite in this section is the fix)
\end{enumerate}
\end{previouslybox}

\begin{topicintro}
Chapter~1 grew its definitions on binary trees, where every node has two named children
and the recurrences close in two lines. Contest trees are not like that: the input is an
edge list, any node may have $10^{5}$ children, and ``left'' and ``right'' do not exist.
This section performs the re-framing once, so that no later chapter has to: root the tree
\emph{inside} the DFS with a parent parameter (\chref{c04}), push \emph{depth} down as an
argument, and let \emph{height} and \emph{size} bubble up as return values -- all three
numbers for all $n$ nodes produced by one traversal (\thmref{c09-onepass}). Nothing here
re-proves \thmref{c01-height-recursion}; the max-and-sum induction is identical, the
children set is merely arbitrary now, and $\mathit{nil}$ has been replaced by the empty
loop.

Two pieces of machinery leave this section and are used book-wide. The first is
\coderef{c09-dfs3} itself -- the three-array DFS that Chapters~13--16 will re-run with
different payloads. The second is \coderef{c09-iter}, the \emph{explicit-stack} rewrite:
when $n$ reaches $2\cdot10^{5}$ and a chain is a legal input, recursion depth is not an
accident but an adversary, and Chapter~16's kit cites \coderef{c09-iter} verbatim for
exactly that case. We finish on LeetCode~742 (\probref{c09-deep}) -- a problem whose whole
trick is running these measures on edges that were never called ``parent'' and ``child''.
\textbf{Prerequisites.} DFS on adjacency lists and the parent-parameter rooting of
Chapter~4 (\chref{c04}); \thmref{c01-height-recursion} as the special case we generalise;
a queue for the BFS inside the solved problem.

\textbf{After this section you can}
\begin{itemize}[nosep]
  \item write the one-DFS triple $(\depth,\height,\size)$ for a general tree cold, and
        state why one traversal suffices; \hfill\lv{2}
  \item rewrite it with an explicit stack and explain precisely what the two-phase order
        buys; \hfill\lv{3}
  \item turn a binary tree into an undirected one and harvest a leaf by distance -- the
        skeleton of \probref{c09-deep}; \hfill\lv{3}
\end{itemize}
\end{topicintro}

\subsection{Intuition}

A rooted tree is a \emph{direction} imposed on a graph. Once you name a root $r$, every
edge acquires an attitude: away from $r$ is down. The DFS discovers that direction by
never stepping back where it came from -- the parent guard is the whole of rooting, no
separate pass needed.

The three measures split by the direction they travel. Depth is inherited: a node learns
its own depth the moment its parent is known ($\depth(\text{child})=\depth(u)+1$), so the
number rides \emph{down} the call as a parameter. Height and size are earned: a node can
only conclude after \emph{every} child has concluded, because $\height(u)$ is a $\max$ and
$\size(u)$ a $\sum$ over the whole child set. One DFS carries both currents at once --
push on the way in, fold on the way back -- which is why three arrays cost no more time
than one.

On a binary tree the fold reads $\max(h_L,h_R)+1$; on a general tree it is a loop, and
the loop is the honest generalisation: $1+\max$ over a set of children of any size.
Chapter~1's conventions $\height(\nil)=-1$ and $\size(\nil)=0$ survive as the
\emph{initialisers} $\height(u)=0$ for a leaf (empty $\max$ leaves it) and $\size(u)=1$
(you count yourself) -- the special case that ``needs no code'' becomes ``a value set
before the loop''.

\subsection{Definitions and Notation}

\begin{definition}{Depth, height and size on a general tree}{c09-measures}
Fix a root $r$ in a tree $T=(V,E)$, $|V|=n$, stored as adjacency lists
$\mathit{adj}[\cdot]$, and fix the \emph{rooted view at $r$} --- every edge oriented
away from $r$. The
\emph{parent} of $u\ne r$ is its neighbour on the path to $r$, its \emph{children} are its
other neighbours, and $\children(u)$ denotes that set ($\children(u)=\emptyset$ makes $u$
a \emph{leaf}). Then
\[
  \depth(u) \;=\; d(r,u),\qquad
  \height(u) \;=\; \max_{v\in\subtree(u)} d(u,v),\qquad
  \size(u) \;=\; |\subtree(u)|,
\]
with $\height(u)=0$ at a leaf and $\size(u)\ge 1$ always ($u$ counts itself). The height
of the tree is $\height(T)=\height(r)$; the root's subtree is all of $V$, so $\size(r)=n$
-- a built-in sanity check.
\end{definition}

\noindent\textbf{In plain words.} Depth looks up (how far from the boss), height and size
look down (how deep my department, how many people in it). Same three questions Chapter~1
asked (\defref{c01-depth}, \defref{c01-height}, \defref{c01-size}); now ``child'' means
\emph{every neighbour except the one I came from}, and that is the entire upgrade.

\begin{theorem}{One DFS computes all three measures in $\bigO{n}$}{c09-onepass}
Algorithm \coderef{c09-dfs3} fills $\depth[\cdot]$, $\height[\cdot]$, $\size[\cdot]$ for
every node of an $n$-node tree in $\Theta(n)$ time and $\bigO{h}$ auxiliary stack space,
and its outputs satisfy \defref{c09-measures}.
\end{theorem}
\begin{proofbox}
The DFS visits each node once (a tree has a unique path from $r$ to every node, so no
node is reached twice) and each edge twice, once per endpoint, of which exactly one
endpoint's loop iteration is rejected by the parent guard. Total loop work is
$\sum_u \deg(u) = 2(n-1) = \Theta(n)$. Correctness is induction on subtree height: when
$\mathit{dfs3}(u,p)$ returns, every child call has returned, so by induction each child
$v$ carries correct $\size(v)$, $\height(v)$; the loop sets
$\size(u)=1+\sum_v\size(v)$ and $\height(u)=\max_v(\height(v)+1)$, matching
\thmref{c01-height-recursion} with the empty-max convention (leaf stays at $0$,
size stays at $1$). Depth is set on entry from the parent's already-final value along the
unique $r$--$u$ path, so $\depth(u)=d(r,u)$. Stack depth is the longest root-to-node
chain, i.e.\ $\bigO{h}$.
\end{proofbox}

\begin{insightbox}
Three \emph{currents} live in one recursion: depth flows down (argument), height and size
flow up (return values). Any quantity in the down-current can be recomputed at any node
on sight; anything in the up-current only becomes final as the call unwinds. That
asymmetry is exactly why the iterative rewrite (\coderef{c09-iter}) needs \emph{two}
phases -- a preorder to plant the down-current, and a reversed replay to finish the
up-current.
\end{insightbox}

\subsection{The Algorithm: One Recursion, Three Arrays}

\begin{lstlisting}[caption={The general-tree measures DFS: depth pushed down, height and size folded up}, label={lst:c09-dfs3}]
// adj[u] = neighbours of u in the undirected tree.  Call dfs3(r, -1).
vector<vector<int>> adj;
vector<int> dep, hgt, sz;                 // sized n

void dfs3(int u, int p) {
    dep[u] = (p == -1 ? 0 : dep[p] + 1);  // depth rides in from the parent
    sz[u]  = 1;   hgt[u] = 0;             // count u; leaf height 0 (edges)
    for (int v : adj[u]) if (v != p) {    // parent guard = the rooting
        dfs3(v, u);                       // child subtree finishes first...
        sz[u]  += sz[v];                  // ...then contributes: sum  -> size
        hgt[u]  = max(hgt[u], hgt[v] + 1);// ...and: max+1 -> height
    }
}
// after the call: tree height = hgt[r], and sz[r] == n is the sanity check
\end{lstlisting}

\begin{complexitybox}
Time $\Theta(n)$ by \thmref{c09-onepass} -- every neighbour slot is inspected exactly
once across the whole run. Memory: three $\bigO{n}$ arrays plus $\bigO{h}$ call stack.
Verdict: the cheapest possible engine for the most asked-for numbers in tree problems;
when the problem asks for \emph{any} of depth, height, size, assume all three came free.
\end{complexitybox}

\subsection{The Chain That Breaks the Stack: \coderef{c09-iter}}

\thmref{c09-onepass}'s honest fine print is the stack term: on a chain of $n$ nodes the
recursion is $n$ frames deep. Constraints of $n \le 2\cdot 10^{5}$ with unrestricted edges
(CSES, most Codeforces) \emph{allow} chains, and contest stacks are small -- the failure
is a runtime error, not a wrong answer, and it will not show on hand-made samples. The
rewrite keeps one plain vector as the stack and splits the visit into two phases:

\begin{lstlisting}[caption={Iterative DFS: preorder to plant depth, reversed replay to finish height and size (the rewrite Chapter 16's kit cites for $5\cdot 10^{5}$-node chains)}, label={lst:c09-iter}]
void dfs3_iter(int r, int n) {
    vector<int> par(n, -1), dep(n), hgt(n), sz(n), order;
    order.reserve(n);
    vector<int> st{r};                       // explicit stack: just nodes
    while (!st.empty()) {                    // phase 1: preorder
        int u = st.back();  st.pop_back();
        order.push_back(u);                  // parent lands in `order` before kids
        for (int v : adj[u]) if (v != par[u]) {
            par[v] = u;  dep[v] = dep[u] + 1;
            st.push_back(v);
        }
    }
    for (int i = (int)order.size() - 1; i >= 0; --i) {   // phase 2: children first
        int u = order[i];                    // (reversed preorder = safe fold order)
        sz[u] = 1;  hgt[u] = 0;
        for (int v : adj[u]) if (v != par[u]) {
            sz[u] += sz[v];
            hgt[u] = max(hgt[u], hgt[v] + 1);
        }
    }
}
\end{lstlisting}

\begin{notebox}[Why the reversed replay is legal]
In a tree every node is pushed only by its parent, so phase~1 writes $\mathit{order}$
parent-before-child. Scanning $\mathit{order}$ backwards therefore visits every child
before its parent -- precisely the dependency the folds in phase~2 need. No
enter/exit-frame bookkeeping, no iterator stack: two passes, one extra vector. (This
trick exists \emph{because} the graph is a tree -- on a general graph a preorder is not a
topological order.) \coderef{c09-iter} replaces \coderef{c09-dfs3} verbatim in every
later chapter's budget analysis; Chapter~16 (\chref{c11}) cites it in its kit for the
same reason.
\end{notebox}

\subsection{Worked Example}

\begin{example}{Rooting the running tree $T$ at node 1: the three arrays}{c09-measure-run}
$T$ has edges $1\!-\!2,1\!-\!3,2\!-\!4,2\!-\!5,3\!-\!6,3\!-\!7,5\!-\!8,7\!-\!9,9\!-\!10$
(\figref{c09-runtree}). Run \coderef{c09-dfs3} with $r=1$; the tables are hand-computed
and fixed for the whole chapter:
\end{example}

\begin{figure}[H]
\centering
\begin{tikzpicture}[diagram, scale=0.85, transform shape,
    tnode/.style={circle, draw=primaryblue, line width=0.9pt, fill=primaryblue!8,
                  minimum size=6.2mm, font=\small\bfseries, text=primaryblue!60!black}]
  \node[tnode] (n1) at (0,0) {$1$};
  \node[tnode] (n2) at (-3,-1.6) {$2$};
  \node[tnode] (n3) at ( 3,-1.6) {$3$};
  \node[tnode] (n4) at (-5.2,-3.2) {$4$};
  \node[tnode] (n5) at (-1.6,-3.2) {$5$};
  \node[tnode] (n6) at ( 1.6,-3.2) {$6$};
  \node[tnode] (n7) at ( 5.2,-3.2) {$7$};
  \node[tnode] (n8) at (-1.6,-4.8) {$8$};
  \node[tnode] (n9) at ( 5.2,-4.8) {$9$};
  \node[tnode] (n10) at ( 5.2,-6.4) {$10$};
  \draw[tedge] (n1)--(n2) (n1)--(n3) (n2)--(n4) (n2)--(n5)
               (n3)--(n6) (n3)--(n7) (n5)--(n8) (n7)--(n9) (n9)--(n10);
  \foreach \nd/\tr in {n1/0:4:10, n2/1:2:4, n3/1:3:5, n4/2:0:1, n5/2:1:2,
                       n6/2:0:1, n7/2:2:3, n8/3:0:1, n9/3:1:2, n10/4:0:1} {
    \node[font=\tiny\ttfamily, text=warningred!75!black, anchor=north]
      at (\nd.south) {\tr};}
  \node[annot, font=\scriptsize, anchor=south west, text=successgreen!50!black]
    at (-0.4,0.55) {root 1: call dfs3(1, -1)};
\end{tikzpicture}
\caption{Every node of $T$ after one DFS, tagged $\depth{:}\height{:}\size$ in red
below its number. Read the two currents: down a column the $\depth$ values increase by
$1$; up a column $\height$ obeys $\max+1$ and $\size$ obeys a sum, e.g.\ node~$5$ has
$\height{=}1$ from child~$8$, and node~$2$ has $\size{=}1+1+2=4$. $\size(1)=10=n$: the
root's count checks the whole run.}
\label{fig:c09-triples}
\end{figure}

\begin{center}
\small
\renewcommand{\arraystretch}{1.15}
\begin{tabular}{l|cccccccccc}
\toprule
node $u$          & 1 & 2 & 3 & 4 & 5 & 6 & 7 & 8 & 9 & 10 \\
\midrule
$\depth(u)$       & 0 & 1 & 1 & 2 & 2 & 2 & 2 & 3 & 3 & 4 \\
$\height(u)$      & 4 & 2 & 3 & 0 & 1 & 0 & 2 & 0 & 1 & 0 \\
$\size(u)$        & 10& 4 & 5 & 1 & 2 & 1 & 3 & 1 & 2 & 1 \\
\bottomrule
\end{tabular}
\end{center}

\begin{tracebox}[Fold order through \coderef{c09-iter} on $T$, root $1$]
phase 1 (preorder, one legal stack order): order = 1, 3, 7, 9, 10, 6, 2, 5, 8, 4
(each child got par[] and dep[] when its parent was popped; e.g. popping 3 sets
par[7]=3, par[6]=3, dep = 2).\\
phase 2 backwards: 4: sz=1 h=0. \;\; 8: sz=1 h=0. \;\; 5: child 8 -> sz=2, h=1.\\
2: children 4,5 -> sz=1+1+2=4, h=max(0+1,1+1)=2. \;\; 10: sz=1 h=0. \;\;
9: child 10 -> sz=2, h=1. \;\; 7: child 9 -> sz=3, h=2. \;\; 6: sz=1 h=0.\\
3: children 6,7 -> sz=1+1+3=5, h=max(0+1,2+1)=3. \;\;
1: children 2,3 -> sz=1+4+5=10 \checkmark (=n), h=max(2+1,3+1)=4. \;\;\checkmark
\end{tracebox}

\begin{insightbox}[Key takeaways]
\begin{enumerate}[nosep, leftmargin=1.6em]
  \item Rooting is a \emph{parameter}, not a data-structure change: $\mathit{adj}$ stays
        undirected; the parent guard orients it on the fly (\defref{c09-measures}).
  \item One DFS, three arrays, $\Theta(n)$ -- depth pushed down, height and size folded
        up (\coderef{c09-dfs3}, \thmref{c09-onepass}).
  \item When chains are legal inputs at $n \ge 2\cdot 10^{5}$, ship the two-phase
        explicit-stack form (\coderef{c09-iter}); reversed preorder is a safe fold order
        \emph{only} because the graph is a tree.
\end{enumerate}
\end{insightbox}

\subsection{Variations and Twists}

\begin{itemize}[nosep]
  \item \textbf{Deepest node(s)}: keep a running $\max$ of $\depth(u)$ inside
        \coderef{c09-dfs3} and collect the argmax -- one line, and it is the whole of the
        ``find the bottom'' family. \hfill\lv{2}
  \item \textbf{Levels and buckets}: $\depth[\cdot]$ \emph{is} the bucket index; any
        per-level aggregation (level sums, level widths) is one array write per node. The
        Euler-tour chapter (\chref{c24}) runs this bucket all the way to path problems.
        \hfill\lv{3}
  \item \textbf{Every root at once}: the tables above are for $r=1$; asking $\size$ under
        \emph{all} roots is the rerooting question -- Chapter~15 (\chref{c20}) owns it;
        \secref{c09-radius} previews the distance-sum version. \hfill\lv{4}
  \item \textbf{Height convention check}: LC~104-style problems count \emph{nodes} on the
        deepest path, our tables count \emph{edges}; the two differ by exactly $1$. Pick
        one convention per problem and check it on a one-node tree before writing the
        fold. \hfill\lv{3}
\end{itemize}

\subsection{Common Pitfalls}

\begin{warningbox}[Trap 1 -- dropping the parent guard]
On an undirected list, \code{for (int v : adj[u]) dfs3(v, u);} without
\code{if (v != p)} walks straight back into the caller: every edge ping-pongs forever
and the run never terminates (or revisits and corrupts $\size$ on a graph with a cycle).
The guard is not an optimisation -- it \emph{is} the rooting (\defref{c09-measures}).
Symptom to recognise: correct answers on stars, stack overflow or wrong sizes on any
deeper tree.
\end{warningbox}

\begin{warningbox}[Trap 2 -- mixing the two height conventions]
$\height(\text{leaf})=0$ (edges) and $\text{depth}(\text{root})=1$ (nodes, LC~104's
habit) in the same program silently shifts every answer by~$1$, and the shift is
invisible on hand samples like a single edge. Pin the convention with a one-node check
(\figref{c09-triples}: $\height(1)=4$, node-count would be $5$), and note that
$\height(\nil)=-1$ from \thmref{c01-height-recursion} is the same bookkeeping written
the other way.
\end{warningbox}

\begin{warningbox}[Trap 3 -- trusting the recursion budget]
Runtime error on a hidden test that is a chain, AC on every sample: this is
\thmref{c09-onepass}'s $\bigO{h}$ stack term being charged in full. At
$n \ge 2\cdot 10^{5}$ with unconstrained edges, prefer \coderef{c09-iter} from the start
-- do not wait for the verdict to tell you. \hfill(\emph{WA clinic}, station~4
of the chapter end walks this exact misfire.)
\end{warningbox}

\subsection{Problem Lab}

The section's anchor is LeetCode~742, \probref{c09-deep}: start at a given node \emph{in
the middle} of a binary tree and find the nearest leaf, where ``nearest'' may walk
\emph{upward} through parents. The binary-tree child pointers only reach down; the first
move is to stop taking the rooted view for granted -- exactly the re-framing of this
section, run in reverse.

% ---------------------------------------------------------------------------
\solvedproblem{c09-deep}{Closest Leaf in a Binary Tree}{LeetCode 742}
  {https://leetcode.com/problems/closest-leaf-in-a-binary-tree/}{LC Medium · L3}
  {general-tree reframing, undirected adjacency, BFS by layers, subtree/leaf flags}

\begin{problemstatement}[Closest Leaf in a Binary Tree (premium-era; statement and samples verified against the live problem)]
Given the \texttt{root} of a binary tree with $n$ nodes labelled $1..\text{}n$ (each
value unique) and a node \texttt{k} with value \texttt{K}, return the value of the
\term{closest leaf} --- a leaf reachable from node \texttt{K} by the minimum number of
moves, where a move crosses one edge in \emph{either} direction (to a child or to the
parent). A leaf means a node with no children in the original \emph{binary} tree; a
single-node tree has its root as a leaf. Constraints: $1 \le n \le 1000$.
\textbf{Example~1}: tree \texttt{[1,3,2]} (node 1 has left 3, right 2), $K=1$
$\Rightarrow$ output \texttt{2} (both children are leaves at distance~$1$; the reference
answer is the smaller label --- our tie-break below matches it). \textbf{Example~2}:
\texttt{[1]}, $K=1$ $\Rightarrow$ \texttt{1}. \textbf{Example~3}: LC's serialisation
\texttt{[1,2,3,4,null,null,null,5,null,6]} (edges $1\!-\!2$, $1\!-\!3$, $2\!-\!4$,
$4\!-\!5$, $5\!-\!6$), $K=2$ $\Rightarrow$ \texttt{3} (distance~$2$ via~$1$ beats
distance~$3$ via~$5,6$).
\end{problemstatement}

\begin{attemptgate}[Attempt it yourself first -- \lv{3}, give yourself 20 minutes]
The trap is treating \code{TreeNode*} as if it were a rooted data structure you must
respect. You are allowed to move upward; say so in data terms before coding.
\peek{From a node, which neighbours can a shortest walk to a leaf use?}
\nudge{All of them: parent and both children are distance-1 moves. That is just the
undirected adjacency list of \secref{c09-measures}, built from child pointers.}
\fullhint{Pass~1: DFS the binary tree, adding an undirected edge between every node and
its parent, and mark the original leaves (no left, no right). Pass~2: BFS from \texttt{K}
on that graph; the first popped node that is marked leaf is the answer --- BFS pops by
nondecreasing distance, so ``first'' $=$ ``closest''. If \texttt{K} is itself a leaf it is
returned at pop~1 (Example~2).}
\end{attemptgate}

\solstep{Understanding the Problem}
Two things must be reconciled. (i)~``Leaf'' is a \emph{rooted} concept --- no children in
the binary tree --- while (ii)~``distance'' is measured on the \emph{unrooted} graph,
where climbing to a parent is legal. A pure down-DFS from \texttt{K} misses leaves above
\texttt{K} entirely (Example~3's answer $3$ lives through the parent). So the input's
rooting is an accident of storage, not of the question: rebuild \defref{c09-measures}'s
view with \texttt{K}'s component structure exposed, then minimise distance.

\solstep{Key Observations}
\begin{observation}
Child pointers give the edges $u\!\to\!u_\text{left},u\!\to\!u_\text{right}$; adding the
reverse edge at every node turns the tree into exactly the adjacency list this section's
\coderef{c09-dfs3} expects --- with $n \le 1000$ any layout works.
\end{observation}
\begin{observation}
On an unweighted graph, BFS pops vertices in nondecreasing distance from the source; the
first \emph{leaf} popped is therefore a closest leaf, and ties are settled by the order
equal-depth vertices are queued (we queue neighbours in ascending value, which reproduces
Example~1's reference answer).
\end{observation}
\begin{observation}
Leaf-ness must be decided from the \emph{original} binary structure (no children), not
from the undirected graph (where degree~$1$ would mislabel a one-child node as leaf and
the single-node tree has degree~$0$ \emph{and} no children --- \texttt{[1]} works either
way, but one-child cases do not).
\end{observation}

\solstep{Think It Through With Me}
\begin{thinkbox}
\deadend{DFS from \texttt{K} downward only, take the shallowest leaf found.}
{Misses every leaf reachable only by going up first --- fails Example~3 ($K=2$: the leaf
$3$ is \emph{above} node~2). Also mixes ``depth below K'' with true distance when K's own
subtree is lopsided.}
\deadend{For every leaf, compute its distance to \texttt{K} by finding both root-paths
and subtracting.}
{Per-leaf distance via a root path is $\bigO{h}$ each after $\bigO{n}$ setup and needs
the LCA path arithmetic of Chapter~16 (\chref{c11}) --- heavyweight for one source, and
error-prone about the common-ancestor formula. We have \emph{one} source node: a single
sweep beats $n$ queries.}
\aha{Undirected edges $+$ one BFS from \texttt{K}. $\bigO{n}$ build, $\bigO{n}$ sweep,
first leaf popped wins --- this section's reframing (\defref{c09-measures}) plus the
queue, nothing more.}
\end{thinkbox}

\solstep{Solution}
Two passes. Pass~1 walks the binary tree and (a) records every original leaf's value and
(b) inserts each tree edge in both directions into $\mathit{adj}[\cdot]$, keeping
neighbour lists ascending so ties resolve to the smaller label. Pass~2 runs BFS from
\texttt{K}: pop, return immediately if the node is a recorded leaf, else queue unseen
neighbours. Correctness rests on one invariant: BFS queues nothing before all
distance-$d$ vertices have been queued, so the first leaf popped is at minimum distance.

\solstep{Proof of Correctness}
\begin{proofbox}[The returned value is a closest leaf of the original tree]
Pass~1 copies every binary-tree edge to the undirected list in both directions, so the
graph is exactly $T$ re-rooted away (\defref{c09-measures}); a walk in the problem's
sense is a path in this graph, of equal length. BFS on any graph pops vertices in order
of nondecreasing distance from the source, since each enqueued vertex sits at
$\mathit{dist}+1$ of its popper and the queue is FIFO: at the time the first leaf $L$ is
popped at distance $d$, no vertex at distance $<d$ remains unpopped, so no leaf is
closer to \texttt{K} than $d$. Leaf membership is tested against the values recorded in
pass~1 from the original child structure, so ``leaf'' matches the problem's definition
including the single-node root (\texttt{[1]} $\Rightarrow$ leaf). A closest leaf always
exists (the finite tree has at least one leaf), so the loop returns before emptying; the
trailing \code{return -1} is unreachable. \texttt{K} itself is popped first at $d=0$,
covering \texttt{K} being a leaf --- Example~2.
\end{proofbox}

\solstep{Implementation}
\begin{lstlisting}[caption={LeetCode 742 -- undirected rebuild, then one BFS by layers}, label={lst:c09-742-soln}]
class Solution {
    set<int> leaves;                       // values with no children, original view
    map<int, set<int>> g;                  // undirected adj; ascending order = tie rule
    void build(TreeNode* u, TreeNode* p) {
        if (!u->left && !u->right) leaves.insert(u->val);      // binary leaf
        if (p) { g[u->val].insert(p->val); g[p->val].insert(u->val); }
        if (u->left)  build(u->left,  u);
        if (u->right) build(u->right, u);
    }
public:
    int findClosestLeaf(TreeNode* root, int K) {
        build(root, nullptr);
        queue<int> q;  set<int> seen;
        q.push(K);  seen.insert(K);
        while (!q.empty()) {
            int u = q.front();  q.pop();
            if (leaves.count(u)) return u; // first leaf popped = closest leaf
            for (int v : g[u]) if (!seen.count(v)) {
                seen.insert(v);  q.push(v);
            }
        }
        return -1;                         // unreachable: every tree has a leaf
    }
};
\end{lstlisting}

\solstep{Trace}
\begin{tracebox}[Official Example~3: edges 1-2, 1-3, 2-4, 4-5, 5-6; K = 2 $\Rightarrow$ 3]
build: leaves = {3, 6} (3 has no children; 6 none). adj: 1->\{2,3\}, 2->\{1,4\},
4->\{2,5\}, 5->\{4,6\}, 6->\{5\}.\\
BFS from 2: pop 2 (not leaf) -> queue 1, 4 (ascending).  pop 1 -> queue 3.
pop 4 -> queue 5.  pop 3 -> LEAF, return \textbf{3}. \checkmark\ matches official.\\
Distances: d(2,3) = 2 < d(2,6) = 3, so 3 is the unique closest leaf.\\
Example 1 (tree 1 with children 3,2; K=1): pop 1 -> queue 2,3 (ascending); pop 2 ->
leaf, return \textbf{2} \checkmark\ (tie at distance 1 broken to the smaller label, as
the reference answer requires). Example 2 (\texttt{[1]}): pop 1 is itself a leaf ->
return \textbf{1} \checkmark.
\end{tracebox}

\solstep{Complexity}
\begin{complexitybox}
Time: $\bigO{n}$ for \code{build} (each edge inserted twice) and $\bigO{n}$ for the BFS
(each vertex queued once) --- the \coderef{c09-dfs3} budget again, with $\log$ factors
from the ordered sets only ($n\le 1000$, irrelevant; swap \code{set} for
\code{unordered\_set} if a tie-break rule is not wanted). Memory: $\bigO{n}$.
Verdict: \ac --- and note what was \emph{not} needed: no LCA, no per-pair distance, just
the reframing of this section plus a queue.
\end{complexitybox}

\solstep{Pitfalls and Variants}
\begin{warningbox}
The two ways to lose this problem: (i) test leaf-ness as ``degree~$1$ in the undirected
graph'' --- a one-child internal node then masquerades as a leaf; (ii) forget the
\texttt{K}$\to$parent escape and DFS only downward (Example~3 wrong). \emph{Variant
worth doing:} LC~865's ``deepest leaves'' family asks the same harvest question but
downward only --- once you have solved \probref{c09-deep} it is the same skeleton with
the up-edges deleted, which is why Chapter~16 cites \probref{c09-deep}, not LCA, for that
shape.
\end{warningbox}

\subsection{Practice Problems}

\begin{stuckbox}
\begin{itemize}[nosep, leftmargin=1.4em]
  \item Did pass~1 record leaf-ness from the \emph{child} structure and adjacency from the
        \emph{undirected} one? They are different facts.
  \item Does your BFS test ``leaf'' at \emph{pop} time, not push time? (A queued vertex
        can be re-discovered earlier; pop time is when distance is final.)
  \item If \texttt{K} is itself a leaf, does your code still return \texttt{K}?
\end{itemize}
\end{stuckbox}

\begin{practiceset}
  \practice{Evaluate Boolean Binary Tree}{LeetCode 2331}
    {https://leetcode.com/problems/evaluate-boolean-binary-tree/}{LC Easy · L2}
    {Leaves carry \texttt{0/1} (false/true); internal nodes carry \texttt{2} (OR) or
    \texttt{3} (AND) --- read the encoding off the statement carefully, it is the other
    way round from what people guess; evaluate the root. This is \coderef{c09-dfs3}'s
    \emph{return-value}
    habit with a different fold: instead of $\max$/$\sum$, the node combines its children
    with the boolean operator it stores --- \texttt{2} combines its two children with
    OR, \texttt{3} with AND, and a leaf just reports its own truth value.
    Check the convention on \texttt{[2,1,3,null,null,0,1]}: node 3 is
    $\mathrm{AND}(0,1)=0$, root is $\mathrm{OR}(1,0)=1$, output \texttt{true}; a
    single \texttt{[0]} root outputs \texttt{false}.}
  \practice{Even Odd Tree}{LeetCode 1609}
    {https://leetcode.com/problems/even-odd-tree/}{LC Medium · L3}
    {Return \texttt{true} iff every even-indexed level holds \emph{odd} values in
    strictly \emph{increasing} order left-to-right and every odd-indexed level holds
    \emph{even} values in strictly \emph{decreasing} order (root at level~0). The level
    index is \coderef{c09-dfs3}'s $\depth[\cdot]$ --- so this is a \emph{per-level}
    check: BFS level by level (queue size per round), and inside each round compare
    neighbours pairwise against the two rules. Official samples:
    \texttt{[1,10,4,3,null,7,9,12,8,6,null,null,2]} $\Rightarrow$ \texttt{true} (levels
    $[1],[10,4],[3,7,9],[12,8,6,2]$); \texttt{[5,4,2,3,3,7]} $\Rightarrow$
    \texttt{false} (level~2 not increasing); \texttt{[5,9,1,3,5,7]} $\Rightarrow$
    \texttt{false} (level~1 not even). $n \le 10^{5}$: mind the chain, BFS is iterative
    anyway.}
\end{practiceset}

\begin{recallbox}[10-second refresher -- 12.1 Depth, Height and Subtree Size]
Root a general tree inside the DFS: parent guard orients the adjacency list
(\defref{c09-measures}). One DFS pushes $\depth$ down and folds $\height$
($\max+1$) and $\size$ (sum) up in $\Theta(n)$ (\coderef{c09-dfs3},
\thmref{c09-onepass}) --- Ch.~1's binary recursions, now over $\children(u)$ of any
size. Chains at $n\ge 2\cdot 10^{5}$? Use the two-phase explicit stack: preorder plants
depth, reversed preorder finishes the folds (\coderef{c09-iter}). \probref{c09-deep}
(LC~742) is the reframing in reverse: undirected rebuild, BFS from \texttt{K}, first
leaf popped wins.
\selftest{
  \item Which of the three measures travel down the recursion, and which travel up?
  \item State the two conventions that make a leaf need no special case in the fold.
  \item Why is reversed preorder a legal fold order on a tree (and what breaks it on a
        general graph)?
  \item In LC~742, why must leaf-ness be tested against the original child structure?
  \item What symptom says ``you should have shipped \coderef{c09-iter}''?
}
\end{recallbox}
